\documentclass[12pt]{article}

\usepackage[T1]{fontenc}
\usepackage[utf8]{inputenc}
\usepackage[brazil]{babel}
\usepackage{sbc-template}
\usepackage{amsmath,graphicx,tikz,url}
\usepackage{microtype}
\usepackage{breakcites}
\usetikzlibrary{arrows.meta}

\title{Avaliação de busca lexical, semântica e híbrida\\
  com reranking e geração aumentada por recuperação}

\author{Julio Cesar Alves Araujo Fuganti\inst{1} (Nº USP 15638792)\\
  Luiz Gabriel Correia dos Santos\inst{1} (Nº USP 15639682)}

% Os e-mails vieram do histórico Git. Confirmar os contatos e a filiação antes da entrega.
\address{Instituto de Ciências Matemáticas e de Computação (ICMC)\\
  Universidade de São Paulo (USP) -- São Carlos, SP, Brasil
  \email{juliocesarfuganti@gmail.com, luizgcorreiadosantos@usp.br}}

\begin{document}

\maketitle

\section{Introdução}

Sistemas de recuperação de informação ordenam documentos de acordo com sua
relevância para uma consulta. Em coleções científicas da área da saúde, perguntas
em linguagem natural podem expressar o mesmo conceito que um artigo usando
termos diferentes. Também há consultas em que uma palavra específica, como o
nome de um alimento ou de uma substância, ajuda a identificar documentos
relevantes. Essas situações motivam a comparação entre métodos baseados em
correspondência de termos e métodos baseados em representações semânticas.

Este trabalho investiga como a combinação dessas abordagens afeta a recuperação
no NFCorpus, uma coleção de recuperação de informação na área da saúde
\cite{boteva2016,thakur2021}. O sistema usa BM25 como referência lexical e um
bi-encoder para a busca semântica. Os rankings são combinados e os primeiros
candidatos são reordenados por um cross-encoder. Um experimento separado avalia
quanto do desempenho perdido com a introdução de erros ortográficos pode ser
recuperado por um corretor.

O trabalho também investiga o efeito da recuperação sobre as respostas de um
modelo de linguagem. A geração aumentada por recuperação, ou RAG, fornece ao
modelo documentos recuperados para a consulta \cite{lewis2020}. O experimento
compara as respostas para as mesmas perguntas em duas condições, com e sem
acesso aos documentos. A comparação considera tanto a relevância dos documentos
que chegaram ao contexto quanto a qualidade e a rastreabilidade das respostas.

\section{Metodologia}

A Figura~\ref{fig:pipeline} mostra as etapas de recuperação e o caminho até a
geração com contexto. O experimento ortográfico usa variantes das consultas em
uma etapa separada. As etapas do fluxo principal recebem as consultas originais.

\begin{figure}[htbp]
\centering
\resizebox{0.98\linewidth}{!}{%
\begin{tikzpicture}[
  etapa/.style={draw, rounded corners=2pt, align=center, font=\small,
    inner xsep=5pt, inner ysep=4pt, minimum height=10mm},
  fluxo/.style={-{Stealth[length=2mm]}, semithick}
]
  \node[etapa] (consulta) at (0,0) {Consulta};
  \node[etapa, minimum width=3.1cm] (lexical) at (3.0,0.75) {Busca lexical\\BM25};
  \node[etapa, minimum width=3.1cm] (semantica) at (3.0,-0.75) {Busca semântica\\Bi-encoder};
  \node[etapa] (rrf) at (5.45,0) {Fusão\\RRF};
  \node[etapa] (rerank) at (7.6,0) {Reranking\\Cross-encoder};
  \node[etapa] (contexto) at (10.05,0) {Seleção\\Top-5};
  \node[etapa] (llm) at (12.3,0) {LLM\\Resposta};

  \draw[fluxo] (consulta.east) -- ++(4mm,0) |- (lexical.west);
  \draw[fluxo] (consulta.east) -- ++(4mm,0) |- (semantica.west);
  \draw[fluxo] (lexical.east) -- (rrf.155);
  \draw[fluxo] (semantica.east) -- (rrf.205);
  \draw[fluxo] (rrf) -- (rerank);
  \draw[fluxo] (rerank) -- (contexto);
  \draw[fluxo] (contexto) -- (llm);
\end{tikzpicture}%
}
\caption{Fluxo de recuperação e geração na condição com RAG.}
\label{fig:pipeline}
\end{figure}

\subsection{Coleção e pré-processamento}

Utilizamos o NFCorpus distribuído pelo \textit{benchmark} BEIR
\cite{boteva2016,thakur2021}. O carregador do BEIR lê
\texttt{corpus.jsonl}, \texttt{queries.jsonl} e os arquivos de julgamentos de
relevância (\textit{qrels}). A coleção contém 3.633 documentos. Os julgamentos
definem as consultas de cada divisão, com 324 em \texttt{dev} e 323 em
\texttt{test}. O conjunto \texttt{dev} serve ao ajuste dos parâmetros de fusão
e correção ortográfica. O conjunto \texttt{test} é reservado à avaliação final.

Cada documento é representado pela concatenação dos campos \texttt{title} e
\texttt{text}. Na busca lexical, documentos e consultas são convertidos para
minúsculas e tokenizados pela expressão regular \texttt{[a-z0-9]+} do módulo
\texttt{re} do Python. O sistema remove as \textit{stopwords} em inglês
fornecidas pelo scikit-learn e mantém tokens de comprimento mínimo um.

Mantivemos uma configuração lexical fixa ao longo dos experimentos, sem
\textit{stemming} ou lematização, para manter o BM25 como uma referência
simples e concentrar a avaliação nos métodos de recuperação, fusão e
reordenação. Preservar as formas dos termos após a normalização facilita
relacionar a pontuação do BM25 às correspondências encontradas no índice.
Essa escolha pode reduzir a cobertura lexical quando consultas e documentos
usam flexões diferentes de uma mesma palavra.

O bi-encoder e o cross-encoder recebem os textos e aplicam seus próprios
tokenizadores.

\subsection{Índice invertido e busca lexical}

O índice invertido foi implementado pelo grupo com dicionários Python. Cada
termo aponta para uma lista invertida que associa o identificador de cada
documento à frequência do termo. Por exemplo, uma lista
\(\{d_1 \mapsto 2,\ d_2 \mapsto 1\}\) indica duas ocorrências em \(d_1\),
uma em \(d_2\) e frequência documental igual a dois. Também são armazenados o
comprimento de cada documento, o total de documentos e o comprimento médio.
O programa registra o tamanho do vocabulário, o número de entradas nas listas
invertidas, o tempo de construção do índice e o tempo médio de consulta.

Escolhemos o Okapi BM25 \cite{manning2008} como referência lexical porque o
modelo combina a importância do termo na coleção, a saturação da frequência
no documento e a normalização pelo comprimento. Para cada consulta, o
algoritmo percorre as listas invertidas dos seus termos e acumula a pontuação
somente para os documentos encontrados. A implementação calcula
\begin{align*}
  s(q,d) &=
  \sum_{t\in V(q)}
  f(t,q)\,\mathrm{IDF}(t)
  \frac{f(t,d)(k_1+1)}
  {f(t,d)+k_1\bigl(1-b+b\,|d|/\overline{L}\bigr)}, \\
  \mathrm{IDF}(t) &=
    \log\!\left(1+\frac{N-\mathrm{DF}(t)+0{,}5}
                       {\mathrm{DF}(t)+0{,}5}\right).
\end{align*}
Aqui, \(V(q)\) representa os termos distintos de \(q\), \(f(t,x)\) é a
frequência do termo em \(x\), \(|d|\) é o comprimento do documento,
\(\overline{L}\) é o comprimento médio e \(N\) é o número de documentos.
Adotamos \(k_1=1{,}5\) e \(b=0{,}75\). O ranking textual exporta até cem
candidatos. Na implementação atual, a lista é mais curta quando há menos de
cem documentos que contêm pelo menos um termo indexado da consulta.

\subsection{Tolerância a erros ortográficos}

O corretor usa um vocabulário de tokens da coleção antes de
\textit{stemming} ou lematização, com as \textit{stopwords} preservadas.
Para um termo fora do vocabulário, um índice de bigramas de caracteres com
marcadores de borda recupera candidatos. O coeficiente de Jaccard filtra as
alternativas com similaridade inferior a \(0{,}35\). A seleção usa distância
de Damerau--Levenshtein de no máximo duas edições. Em caso de empate, prioriza
maior Jaccard, maior número de documentos com o candidato e ordem alfabética.

O experimento usa dez consultas de \texttt{test} cujo NDCG@10 lexical original
é positivo. Os erros introduzidos abrangem inserção, remoção, substituição e
troca de posição de caracteres. A palavra alterada deve ficar fora do
vocabulário. São comparadas as versões original, com erro e corrigida.
Quando há perda, calcula-se a fração recuperada dividindo a melhora da
versão corrigida sobre a versão com erro pela perda em relação à original.
As mesmas variantes são avaliadas na busca semântica como análise
complementar.

\subsection{Recuperação semântica}

O bi-encoder \texttt{all-MiniLM-L6-v2}, da biblioteca Sentence Transformers,
transforma cada texto em um vetor de 384 dimensões. Os vetores dos documentos
são pré-calculados e armazenados. Consultas e documentos têm vetores
normalizados pela norma \(L_2\). O produto escalar entre eles equivale à
similaridade do cosseno e é calculado com NumPy para todos os documentos.
O tokenizador limita a entrada a 256 tokens. O programa contabiliza os
documentos sujeitos a truncamento e registra o tempo de geração dos
vetores da coleção e o tempo médio por consulta.

\subsection{Busca híbrida}

A busca híbrida usa \textit{Reciprocal Rank Fusion} (RRF)
\cite{cormack2009}. A fusão recebe os primeiros cem documentos disponíveis de
cada ranking. Cada documento \(d\) recebe a pontuação
\[
  \mathrm{RRF}(d)=\sum_{s:\, d\in L_s}\frac{1}{k+r_s(d)},
\]
em que \(L_s\) é a lista de candidatos do sistema \(s\) e \(r_s(d)\) é sua
posição, contada a partir de um. A fusão utiliza posições, o que permite
combinar listas pontuadas em escalas diferentes. Os valores
\(k\in\{10,60,100\}\) são comparados por NDCG@10 médio em \texttt{dev}.
A configuração de maior média é usada na avaliação em \texttt{test}.

\subsection{Reranking}

Utilizamos o cross-encoder \texttt{ms-marco-MiniLM-L6-v2}, pré-treinado para
ordenar passagens no MS MARCO \cite{crossencoder}. Seu treinamento é
adequado à tarefa de pontuar candidatos para uma consulta. O modelo processa
os dois textos conjuntamente e considera interações entre seus tokens.
Esse processamento é limitado aos 50 primeiros candidatos da lista
híbrida para conter o custo por consulta. A ordenação neural mantém a
posição híbrida como critério de desempate e preserva a ordem relativa dos
documentos a partir da posição 51. O modelo recebe até 512 tokens por par.
O programa registra o tempo da reordenação e contabiliza os pares que
ultrapassam o limite.

\subsection{Geração aumentada por recuperação}

O experimento de RAG usa as dez perguntas fixas de \texttt{test}. Seus
identificadores são sorteados com semente 42 de uma lista ordenada de
consultas que terminam em ponto de interrogação. Cada pergunta é respondida em
português pelo \texttt{Qwen/Qwen3-4B}, executado localmente com Transformers,
LangChain e quantização NF4 de quatro bits via bitsandbytes. As condições
compartilham temperatura zero, semente 42, limite de 1.024 tokens de saída e
modo de raciocínio desativado.

Na condição A, o modelo recebe a pergunta sem contexto externo. Na condição B,
o \textit{prompt} inclui a mesma pergunta e os cinco primeiros documentos do
ranking reranqueado. Cada documento é mantido como uma unidade com
identificador, título e texto. A instrução pede uma resposta apoiada no
contexto, com indicação dos identificadores que sustentam as afirmações e
menção explícita quando a evidência é insuficiente. O protocolo prevê avaliar
a relevância dos documentos enviados e a qualidade das respostas por
julgamento humano nas duas condições.

% Próximas seções: Avaliação, Resultados Obtidos e Análise,
% Considerações finais e Uso de ferramentas de IA.
% Ao fechar a Avaliação, registrar as revisões dos pesos, o ambiente e
% os parâmetros da execução escolhida. Atualizar a descrição do Top-100
% lexical se o comportamento do ranking for ajustado.

\begin{thebibliography}{}
\bibitem[Boteva et al. 2016]{boteva2016}
Boteva, V., Gholipour, D., Sokolov, A. e Riezler, S. (2016).
\newblock A Full-Text Learning to Rank Dataset for Medical Information Retrieval.
\newblock In \textit{European Conference on Information Retrieval}.

\bibitem[Cormack et al. 2009]{cormack2009}
Cormack, G. V., Clarke, C. L. A. e Büttcher, S. (2009).
\newblock Reciprocal Rank Fusion Outperforms Condorcet and Individual Rank
Learning Methods.
\newblock In \textit{Proceedings of the 32nd International ACM SIGIR Conference
on Research and Development in Information Retrieval}.

\bibitem[Lewis et al. 2020]{lewis2020}
Lewis, P., Perez, E., Piktus, A. et al. (2020).
\newblock Retrieval-Augmented Generation for Knowledge-Intensive NLP Tasks.
\newblock In \textit{Advances in Neural Information Processing Systems}, vol.~33.

\bibitem[Manning et al. 2008]{manning2008}
Manning, C. D., Raghavan, P. e Schütze, H. (2008).
\newblock \textit{Introduction to Information Retrieval}.
\newblock Cambridge University Press.

\bibitem[Sentence Transformers]{crossencoder}
Sentence Transformers.
\newblock \textit{Cross-Encoder for MS Marco}. Ficha do modelo.
\newblock \url{https://huggingface.co/cross-encoder/ms-marco-MiniLM-L6-v2}.

\bibitem[Thakur et al. 2021]{thakur2021}
Thakur, N., Reimers, N., Rücklé, A., Srivastava, A. e Gurevych, I. (2021).
\newblock BEIR: A Heterogeneous Benchmark for Zero-shot Evaluation of
Information Retrieval Models.
\newblock In \textit{NeurIPS Datasets and Benchmarks}.
\end{thebibliography}

\end{document}
